\chapter{روش تحقیق و طراحی چارچوب آزمایش}

\section{طرح کلی پژوهش}

این پژوهش یک مطالعه تجربی و مقایسه‌ای برای بررسی روش‌های گراف‌محور در تشخیص توضیح‌پذیر سلامت روان است. هدف، مقایسه روش‌ها در یک وظیفه و با منابع داده و دانش مشترک است تا تفاوت مشاهده‌شده تا حد امکان ناشی از سازوکار بازیابی و استدلال هر روش باشد، نه تفاوت در داده یا معیار ارزیابی. چارچوب آزمایش برای دو مجموعه‌داده \lr{DR} و \lr{SWMH} اجرا می‌شود و روش‌های \lr{ToG-1}، \lr{ToG-2}، \lr{HippoRAG 2}، بازیابی افزوده متعارف (\lr{Vanilla RAG}) و استدلال زنجیره‌ای (\lr{CoT}) را در بر می‌گیرد.

ورودی هر آزمایش یک متن از مجموعه‌داده به همراه مجموعه کامل برچسب‌های ممکن آن وظیفه است. متن در صورت نیاز به موجودیت‌ها و مفاهیم مرتبط پیوند داده می‌شود و شواهد دانشی و متنی از گراف دانش و مخزن قطعه‌ها بازیابی می‌گردند. گراف دانش دامنه‌ای از داده‌های مورد استفاده در این پژوهش ساخته شده و سه‌تایی‌ها، موجودیت‌ها و روابط آن در پایگاه داده \lr{PostgreSQL} نگهداری می‌شوند؛ نمایش‌های برداری موجودیت‌ها و قطعه‌های متنی نیز در \lr{Milvus} در دسترس روش‌های بازیابی قرار دارند.

هر روش پس از دریافت ورودی، خروجی خود را در قالب یک پیش‌بینی از میان برچسب‌های مجاز تولید می‌کند. افزون بر برچسب، توضیح متنیِ دلیل تصمیم و شواهد مورد استفاده نیز ثبت می‌شود. در روش‌های گراف‌محور، مسیر یا زیرساختارهای گرافی، موجودیت‌ها، روابط و قطعه‌های انتخاب‌شده ثبت می‌شوند؛ در روش‌های بازیابی‌محور، قطعه‌ها و شواهد بازیابی‌شده؛ و در روش \lr{CoT}، متن استدلال و پاسخ نهایی ذخیره می‌گردد. بنابراین، خروجی آزمایش فقط یک برچسب نیست، بلکه شامل پاسخ قابل ارزیابی و ردپای شواهدی است که به تولید آن پاسخ منجر شده‌اند.

برای هر اجرا، پیش‌بینی با برچسب واقعی مقایسه می‌شود و مراحل داخلی، تصمیم‌های انتخاب و حذف نامزدها، تعداد فراخوانی‌های مدل زبانی، توکن‌های مصرف‌شده و زمان اجرا ثبت می‌گردند. این ثبت مشترک امکان ارزیابی سه جنبه را فراهم می‌کند: درستی پیش‌بینی، قابلیت بازبینی و پشتیبانی توضیح با شواهد، و هزینه محاسباتی و متنی روش. جزئیات آماده‌سازی داده و گراف، تنظیمات هر روش، پروتکل اجرای آزمایش و معیارهای عددی در بخش‌های بعدی این فصل ارائه می‌شوند.

\section{تعریف وظیفه و خروجی‌های مورد انتظار}

وظیفه اصلی پژوهش، نگاشت یک متن مرتبط با سلامت روان به یکی از برچسب‌های از پیش تعیین‌شده و ارائه توضیحی برای این تصمیم است. هر نمونه شامل یک متن یا پست، شناسه نمونه و برچسب واقعی است. متن خام، همراه با دستور مسئله و فهرست کامل کلاس‌های مجاز، به هر روش داده می‌شود تا مدل فقط یکی از برچسب‌های تعریف‌شده برای همان مجموعه‌داده را انتخاب کند. این محدودیت برای جلوگیری از تولید نام کلاس‌های جدید و قابل مقایسه بودن خروجی روش‌ها اعمال می‌شود.

دو مجموعه‌داده از نظر صورت‌بندی مسئله متفاوت‌اند. در مجموعه‌داده \lr{DR}، مسئله دودویی است و خروجی باید دقیقاً یکی از دو کلاس \lr{yes} یا \lr{no} باشد؛ این دو کلاس به‌ترتیب وجود یا نبود وضعیت هدف در نمونه را نشان می‌دهند. در مجموعه‌داده \lr{SWMH}، مسئله چندکلاسه است و مدل باید یکی از پنج کلاس \lr{anxiety}، \lr{bipolar disorder}، \lr{depression}، \lr{no mental disorders} یا \lr{suicide} را انتخاب کند. برچسب واقعی هر نمونه از فایل داده خوانده و پس از اجرای روش با برچسب پیش‌بینی‌شده مقایسه می‌شود.

خروجی هر اجرا فقط به یک برچسب محدود نیست. بخش نخست خروجی، پاسخ نهایی شامل برچسب پیش‌بینی‌شده و توضیح متنی دلیل انتخاب آن است. بخش دوم، شواهدی است که روش برای رسیدن به پاسخ استفاده کرده است. در روش‌های گراف‌محور، این شواهد شامل موجودیت‌های آغازین، روابط و مسیرهای انتخاب‌شده و قطعه‌های متنی متصل به آن‌هاست. در روش‌های بازیابی‌محور، قطعه‌های متنی بازیابی‌شده و در روش \lr{CoT}، متن استدلال و پاسخ تولیدشده ثبت می‌شود. این اطلاعات امکان بررسی این موضوع را فراهم می‌کند که پاسخ بر چه شواهدی تکیه داشته است.

برای هر نمونه، علاوه بر خروجی محتوایی، برچسب واقعی، وضعیت درستی پیش‌بینی و اطلاعات اجرای روش نیز ذخیره می‌شود. اطلاعات اجرایی شامل فراخوانی‌های مدل زبانی، مسیر مراحل بازیابی یا پیمایش، توکن‌های ورودی و خروجی، زمان هر فراخوانی و زمان کل اجراست. در صورتی که مدل پاسخی خارج از قالب مورد انتظار تولید کند یا اجرای روش با خطای زیرساختی متوقف شود، این وضعیت جداگانه ثبت می‌شود تا با یک پیش‌بینی عادی اشتباه نشود.

لازم است تأکید شود که برچسب تولیدشده در این پژوهش معادل تشخیص قطعی بالینی نیست. خروجی، پیش‌بینی محاسباتی مدل بر اساس متن نمونه و شواهد در دسترس است و توضیح تولیدشده نیز باید به عنوان توضیح قابل بازبینیِ تصمیم مدل تفسیر شود، نه جایگزین ارزیابی متخصص. معیار مقایسه در این پژوهش، توافق برچسب با داده مرجع، پشتیبانی توضیح از شواهد و هزینه تولید پاسخ است.

\section{منابع داده و گراف دانش دامنه‌ای}

منبع دانش دامنه‌ای پژوهش از دو مرجع تخصصی حوزه اختلال‌های روانی تهیه شده است: \lr{Diagnostic and Statistical Manual of Mental Disorders, Fifth Edition, Text Revision (DSM-5-TR)} منتشرشده توسط انجمن روان‌پزشکی آمریکا و کتاب \lr{DSM-5-TR Made Easy: The Clinician's Guide to Diagnosis} نوشته جیمز موریسون \cite{apa2022dsm5tr,morrison2022dsm5trmadeeasy}. مرجع نخست چارچوب رسمی طبقه‌بندی و معیارهای تشخیصی را فراهم می‌کند و مرجع دوم همان مفاهیم و معیارها را با بیان آموزشی و ساختار مناسب‌تر برای استخراج روابط توضیح می‌دهد. محتوای این دو منبع پس از استخراج مفاهیم و روابط، در گراف دانش دامنه‌ای یکپارچه شده است.

گراف حاصل، مفاهیمی مانند اختلال‌ها، نشانه‌ها، رفتارها، معیارها و مشخصه‌های مرتبط و روابط میان آن‌ها را نگهداری می‌کند. نمایش اصلی گراف به صورت سه‌تایی‌های \lr{(subject, predicate, object)} است. موجودیت‌ها و نوع آن‌ها، سه‌تایی‌ها و رابطه‌ها، و پیوند میان سه‌تایی‌ها و قطعه‌های منبع در پایگاه داده رابطه‌ای \lr{PostgreSQL} ذخیره شده‌اند. شناسه‌های موجودیت و قطعه در لایه‌های مختلف حفظ می‌شوند تا هر شاهد بازیابی‌شده به رکورد گراف و متن منبع آن قابل ردیابی باشد.

متن‌های منبع به قطعه‌های مستقل تقسیم شده‌اند و برای هر قطعه شناسه و محتوای متنی ثبت شده است. جدول پیوندِ ذکرها مشخص می‌کند هر قطعه از کدام سه‌تایی‌ها پشتیبانی می‌کند. در نتیجه، بازیابی یک موجودیت یا یک مسیر گرافی می‌تواند به قطعه‌های متنی مرتبط برسد و روش‌های \lr{ToG-1} و \lr{ToG-2} قادرند هم ساختار رابطه‌ای و هم زمینه متنی را در اختیار مدل قرار دهند. هنگام حذف نامزدهای تکراری، شناسه قطعه و شناسه سه‌تایی به عنوان شناسه‌های مستقل حفظ می‌شوند؛ بنابراین، یک قطعه مشترک می‌تواند به چند موجودیت یا چند رابطه متصل باشد، بدون آن‌که به‌اشتباه چند قطعه متفاوت تلقی شود.

برای پیوند موجودیت‌های استخراج‌شده از پرسش، از نمایش‌های برداری ذخیره‌شده در \lr{Milvus} استفاده می‌شود. متن موجودیت یا عبارت استخراج‌شده با مدل \lr{embedding} سرویس \lr{Jina} به بردار تبدیل و با جست‌وجوی شباهت کسینوسی با موجودیت‌های نمایه‌شده مقایسه می‌شود. شناسه‌های موجود در مجموعه‌های \lr{Milvus} با شناسه‌های موجودیت در \lr{PostgreSQL} یکسان نگه داشته شده‌اند تا نتیجه جست‌وجوی برداری مستقیماً به گره گراف قابل اتصال باشد. بردارهای قطعه‌های متنی نیز در مجموعه جداگانه‌ای ذخیره شده‌اند و برای رتبه‌بندی شواهد با بردار پرسش به کار می‌روند.

این جداسازی سه لایه دارد: \lr{PostgreSQL} منبع ساختاری گراف و رابطه‌های اتصال‌دهنده است؛ \lr{Milvus} نمایه جست‌وجوی برداری موجودیت‌ها و قطعه‌هاست؛ و سرویس \lr{embedding} تبدیل متن به بردار را انجام می‌دهد. مدل زبانی از طریق \lr{OpenRouter} فقط در مراحل استخراج موجودیت، انتخاب یا ارزیابی نامزدها و تولید پاسخ به کار می‌رود. جزئیات نحوه استفاده هر روش از این منابع در بخش روش‌های مورد مقایسه و پروتکل اجرای آزمایش شرح داده می‌شود.

\section{روش‌های مورد مقایسه}

در این پژوهش پنج روش برای تولید پیش‌بینی و توضیح روی دو مجموعه‌داده با یکدیگر مقایسه می‌شوند. روش‌ها از نظر میزان استفاده از گراف دانش، نحوه بازیابی شواهد و میزان دخالت مدل زبانی در فرایند استدلال متفاوت‌اند، اما همگی در نهایت باید یک کلاس مجاز و توضیح مرتبط با همان نمونه تولید کنند. اجرای روش‌ها با منابع داده و گراف دامنه‌ای یکسان انجام می‌شود تا مقایسه بر تفاوت سازوکار روش‌ها متمرکز باشد.

\subsection{استدلال زنجیره‌ای و بازیابی افزوده متعارف}

روش \lr{CoT} به عنوان خط پایه‌ی زبانی، متن نمونه و دستور طبقه‌بندی را در اختیار مدل قرار می‌دهد و از مدل می‌خواهد ابتدا درباره‌ی شواهد موجود استدلال و سپس یک کلاس نهایی انتخاب کند. در این روش، پیمایش صریح گراف وجود ندارد و شواهد اصلی از متن ورودی و دانش درونی مدل به دست می‌آیند. خروجی شامل متن استدلال، پاسخ نهایی و اطلاعات اجرای فراخوانی مدل است.

در روش \lr{Vanilla RAG}، پیش از تولید پاسخ، قطعه‌های متنی مرتبط با پرسش از نمایه برداری بازیابی می‌شوند. قطعه‌های منتخب همراه با پرسش به مدل داده می‌شوند و مدل با اتکا به این بافت، لیبل و توضیح را تولید می‌کند. این روش نقش گراف را در هدایت چندمرحله‌ای جست‌وجو وارد نمی‌کند و مبنایی برای سنجش اثر افزودن ساختار گرافی به بازیابی فراهم می‌آورد.

\subsection{روش‌های گراف‌محور}

\lr{ToG-1} روش عامل‌محوری است که جست‌وجو را از موجودیت‌های مرتبط با پرسش آغاز می‌کند و در هر مرحله، روابط و موجودیت‌های نامزد را با کمک مدل زبانی پالایش می‌کند. مسیرهای برتر برای گسترش مرحله بعد حفظ می‌شوند و در صورت کافی بودن شواهد، مدل پاسخ نهایی و توضیح آن را تولید می‌کند. در این روش، مسیرهای گرافی نقش اصلی را در سازمان‌دهی شواهد دارند.

\lr{ToG-2} نسخه‌ای توسعه‌یافته از این الگوست که بازیابی گراف و بازیابی متن را در یک چرخه‌ی مشترک قرار می‌دهد. ابتدا موجودیت‌های آغازین از متن استخراج و به گره‌های گراف پیوند داده می‌شوند. سپس قطعه‌های متنی مرتبط با \lr{seed}ها و نامزدهای گراف بر اساس شباهت با پرسش رتبه‌بندی می‌شوند. مدل در هر دور درباره‌ی کافی بودن شواهد تصمیم می‌گیرد؛ در صورت ناکافی بودن، روابط، موجودیت‌ها و \lr{chunk}های دور بعد انتخاب می‌شوند و این فرایند تا رسیدن به پاسخ یا حد عمق ادامه می‌یابد.

\lr{HippoRAG 2} به عنوان روش بازیابی‌محور گرافی، از ساختار گراف و ارتباط میان مفاهیم برای انتشار یا تجمیع شواهد مرتبط استفاده می‌کند. در این روش، گراف نقش حافظه‌ی غیرپارامتری را دارد و بازیابی صرفاً بر شباهت مستقل پرسش با یک قطعه متکی نیست. خروجی بازیابی‌شده به مدل داده می‌شود تا برچسب و توضیح نهایی را تولید کند. بنابراین، این روش در مقایسه با \lr{Vanilla RAG} اثر استفاده از پیوندهای گرافی را نشان می‌دهد، اما مانند \lr{ToG-2} چرخه‌ی تصمیم‌گیری عامل‌محور برای ادامه پیمایش ندارد.

این پنج روش یک طیف از استدلال صرفاً زبانی تا بازیابی متنی، بازیابی گرافی و پیمایش عامل‌محور گراف را پوشش می‌دهند. در تمام روش‌ها، متن پرسش، مجموعه برچسب‌های مجاز و تعریف خروجی یکسان نگه داشته می‌شود؛ تفاوت اصلی در منبع شواهد و نحوه انتخاب و سازمان‌دهی آن‌هاست. جزئیات تنظیمات مدل، پارامترهای بازیابی، بودجه جست‌وجو و ترتیب اجرای هر روش در بخش پروتکل آزمایش ارائه خواهد شد.

\section{چارچوب اجرای آزمایش‌ها}

آزمایش‌ها بر روی نسخه‌ی ثابت مجموعه‌داده‌ی مورد استفاده انجام شدند. مجموعه‌ی آزمون
شامل دو وظیفه‌ی \lr{DR} و \lr{SWMH} است و در اجرای مقایسه‌ای، به‌ترتیب ۷۰ و ۶۵
نمونه از آن‌ها بررسی شد. برای همه‌ی روش‌ها متن ورودی، برچسب واقعی و مجموعه‌ی کامل
برچسب‌های مجاز یکسان بود.

مقایسه میان پنج روش \lr{CoT}، \lr{Vanilla RAG}، \lr{HippoRAG 2}، \lr{ToG-1} و
\lr{ToG-2} انجام شد. مدل زبانی مشترک، \lr{Gemma 4 31B} \cite{gemmateam2026gemma4}
و مدل تولیدکننده‌ی نمایش برداری، \lr{jina-embeddings-v4}
\cite{gunther2025jinaembeddingsv4universalembeddingsmultimodal} بود. گراف دانش از طریق
\lr{PostgreSQL} و نمایه‌ی برداری از طریق \lr{Milvus} در اختیار روش‌های مربوط قرار گرفتند.

در \lr{CoT}، پاسخ مستقیماً بر اساس متن ورودی تولید شد. در \lr{Vanilla RAG}، ابتدا
۵۰ قطعه‌ی متنی بازیابی و پس از رتبه‌بندی، حداکثر ۸ قطعه به مدل داده شد. در
\lr{HippoRAG 2}، سه‌تایی‌های حقیقت از گراف بازیابی و امتیاز آن‌ها با استفاده از
رتبه‌بندی صفحه‌ی شخصی‌سازی‌شده\footnote{\lr{Personalized PageRank (PPR)}} به قطعه‌های متنی منتقل شد. در \lr{ToG-1}، عرض جست‌وجو ۵ و عمق آن ۵ و
در \lr{ToG-2}، عرض جست‌وجو ۵ و عمق آن ۵ در نظر گرفته شد.

برای هر نمونه، علاوه بر برچسب و توضیح، شواهد بازیابی‌شده، مسیرهای گرافی، زمان
اجرا، تعداد فراخوانی‌های مدل و تعداد توکن‌های ورودی و خروجی ثبت شد. پاسخ خارج از
مجموعه‌ی برچسب‌های مجاز یا پاسخ غیرقابل‌تجزیه، معتبر در نظر گرفته نشد. سپس عملکرد
روش‌ها با استفاده از معیارهای درستی پیش‌بینی، کیفیت توضیح و هزینه‌ی اجرا مقایسه
شد. در اجرای \lr{ToG-2}، دو نمونه خروجی کامل نداشتند؛ بنابراین نتایج این روش بر
اساس ۱۳۳ نمونه‌ی معتبر گزارش شد.

\section{معیارهای ارزیابی}
\mbox{}\par

\section{مقایسه منصفانه، بازتولیدپذیری و جمع‌بندی}
\mbox{}\par
